\chapterbackmatter{Weinberg Exercises}

\begin{enumerate}
\WeinbergExercise{1} Consider a neutral vector field $v_\mu(x)$. What conditions must be
imposed on the sum $\Pi^*_{\mu\nu}(k)$ of one-particle-irreducible graphs with
two external vector field lines in order that the field should be properly
renormalized and describe a particle of renormalized mass $m$? How do we split
the free-field and interacting terms in the Lagrangian to achieve this?

\WeinbergExercise{2} Derive the generalized Ward identity that governs the electromagnetic
vertex function of a charged scalar field.

\WeinbergExercise{3} What is the most general form of the matrix element
$\bra{p_2,\sigma_2}J^\mu(x)\ket{p_1,\sigma_1}$ of the electromagnetic current
$J^\mu(x)$ between two spin $\tfrac12$ one-particle states of \emph{different}
masses $m_1$ and $m_2$ and equal parity? What if the parities were opposite?
(Assume parity conservation throughout.)

\WeinbergExercise{4} Derive the spectral (K\"all\'en--Lehmann) representation for the vacuum
expectation value $\bra{\Omega}T\{J^\mu(x)J^\nu(y)\}\ket{\Omega}$, where
$J^\mu(x)$ is a complex conserved current.

\WeinbergExercise{5} Derive the spectral (K\"all\'en--Lehmann) representation for the vacuum
expectation value
$\bra{\Omega}T\{\psi_n(x)\bar\psi_m(y)\}\ket{\Omega}$, where $\psi(x)$ is a
Dirac field.

\WeinbergExercise{6} Without using any assumptions about the asymptotic behavior of the
scattering amplitude or cross-sections, show that it is impossible for forward
photon scattering amplitudes to satisfy unsubtracted dispersion relations.

\WeinbergExercise{7} Derive the spectral (K\"all\'en--Lehmann) representation for a complex
scalar field by using the methods of dispersion theory.

\WeinbergExercise{8} Use dispersion theory and the results of Section 8.7 to calculate the
amplitude for forward photon--electron scattering in the electron rest frame
to order $e^4$.
\end{enumerate}
